% =====================================================================
%  Chapter 6 -- Conclusion
% =====================================================================
\chapter{Conclusion}
\label{ch:conclusion}

As PECs become more power-dense, removing their losses from a smaller volume
increasingly limits their design. This project evaluated MHD-pumped Galinstan as
a coolant for such converters, and identified where in an IPT system it is most
useful. A Galinstan loop and a water loop were compared on the same commercial
cold plate under resistive heat loads, the results were interpreted with an
$\varepsilon$--NTU thermal network and extended with a validated conjugate heat
transfer simulation, and the Galinstan loop was used to cool a power-dense
\SI{85}{\kilo\hertz} planar inductor.

\textbf{RQ1.} At the flows each pump could deliver, water outperformed
Galinstan, with a total resistance from the resistor to the coolant inlet of
\SIrange{0.257}{0.267}{\celsius\per\watt} against
\SIrange{0.488}{0.799}{\celsius\per\watt}. The $\varepsilon$--NTU analysis
explains this result and supports H1. Galinstan's high thermal conductivity
reduced its convective resistance to less than \SI{3}{\percent} of the cold-plate
resistance, but at \SIrange{0.77}{1.45}{\milli\litre\per\second} the plate was
caloric-limited, so Galinstan's low volumetric heat capacity set its
performance. The water plate, in contrast, was convection-limited and barely
responded to a seven-fold change in flow. The MHD pump drew only
\SI{0.85}{\watt}, against \SI{11.5}{\watt} for the water pump, giving a maximum
COP of 200 against 47.6, but its mechanical efficiency was only
\SI{0.56}{\percent}. The simulation shows that at
\SI{5}{\milli\litre\per\second}, the limit of laminar flow, the Galinstan total
resistance would fall to about \SI{0.29}{\celsius\per\watt}, within
\SI{14}{\percent} of water at a lower flow than any of the water tests.

\textbf{RQ2.} The Galinstan loop cooled the inductor, running at more than twice
the loss its core can dissipate without active cooling, to a maximum surface
temperature of about \SI{83}{\celsius}. Galinstan inside the cold plate added
only \SI{1.0}{\milli\ohm} to the winding AC resistance at
\SI{85}{\kilo\hertz}, against \SI{11.7}{\milli\ohm} for the plate itself and
\SI{10.3}{\milli\ohm} for unshielded Galinstan beside the inductor. The metal of
the cold plate therefore shields the coolant from the leakage field, supporting
H2, although the effect of the field on the flow could not be confirmed.

The main contribution of this work is to show when Galinstan is worth using.
Its advantage lies not in its conductivity alone, but in applications where the
heated area is small and the heat flux high: at \SI{5}{\milli\litre\per\second},
the model predicts that it outperforms water once the heated area falls below
about \SI{40}{\centi\metre\squared}. In an IPT system, this points to the
inverter and rectifier modules and to shielded compact magnetics, rather than
the large, distributed losses of the pads. These conclusions rest on a small
number of tests on a single cold plate designed for water, and the area
analysis is indicative only. Realising Galinstan's advantage will require a more
efficient, higher-flow pump with integrated power electronics, and a smaller cold
plate designed for liquid metal (\cref{sec:disc-further-work}).
